\subsection{Advantages and Limitations}\label{sec:advantages-and-limitations}

\subsubsection{Strengths of Variational Algorithms}\label{sec:strengths-of-variational-algorithms}

Now that we have introduced the main components of variational quantum algorithms, we can discuss their strengths and limitations. We move from \textbf{how VQAs work} to \textbf{why they became the dominant NISQ algorithms}.

\highspace
The main strengths of VQAs are:
\begin{enumerate}
    \item \textbf{Shallow circuits are more resilient to noise.}
    \item \textbf{Repeated sampling and optimization make them robust to errors.}
    \item \textbf{They are highly flexible.}
\end{enumerate}
For these reasons, they are considered \hl{the most promising algorithms for NISQ devices.} Let's examine each point.

\highspace
\begin{flushleft}
    \textcolor{Green3}{\faIcon{shield-alt} \textbf{Shallow Circuits Are More Resilient to Noise}}
\end{flushleft}
This is probably the biggest advantage. Recall the main limitation of NISQ hardware: \textbf{more gates} lead to \textbf{more accumulated errors}, which \textbf{lowers} the \textbf{probability of obtaining the correct result}.

\highspace
A traditional quantum algorithm may require thousands of gates, whereas a VQA typically executes a much shorter circuit. Because the \textbf{circuit is shorter}, \textbf{qubits spend less time evolving}, \textbf{fewer gate errors accumulate}, and \textbf{decoherence has less time to affect the computation}. Hence, shallow circuits lead to less noise and \textbf{more reliable execution}.

\highspace
\begin{flushleft}
    \textcolor{Green3}{\faIcon{sync-alt} \textbf{Repeated Sampling and Optimization Improve Robustness}}
\end{flushleft}
At first glance, \hl{repeatedly executing the circuit may seem inefficient.} However, repetition provides an important benefit.

\highspace
Instead of trusting a single noisy measurement, the algorithm runs the circuit many times and \hl{estimates the expectation value from multiple samples.} The optimizer then uses this more reliable estimate. So the repeated execution is not wasted effort, it \textbf{helps reduce the impact of statistical fluctuations and measurement noise}. This is why repeated sampling and optimization make VQAs more robust to errors.

\highspace
\begin{flushleft}
    \textcolor{Green3}{\faIcon{tools} \textbf{Ample Flexibility}}
\end{flushleft}
One of the nicest features of VQAs is that the \textbf{framework is very general}. Remember the three building blocks: \textbf{ansatz}, \textbf{cost function}, and \textbf{optimizer}. Each can be \textbf{modified independently}. This modularity allows researchers to adapt the algorithm to many different applications without redesigning everything from scratch.

\highspace
In summary, VQAs are the most promising algorithms for NISQ devices. Despite the three major limitations of NISQ devices (noise, few qubits, and limited circuit depth, \autopageref{sec:why-variational-algorithms}), \textbf{VQAs address these limitations} by using shallow circuits, which reduce decoherence and gate errors; hybrid optimization, which allows for repeated measurements and better estimation of the cost function; and flexibility, which allows different ansätze, optimizers, and cost functions to be combined for various applications. Rather than trying to eliminate noise completely, VQAs are designed to \textbf{work despite the presence of noise}.